\begin{abstract}
The rapid expansion of Augmented Reality (AR) on the web is currently hindered by the significant computational overhead of existing frameworks, which typically rely on heavy machine learning libraries such as TensorFlow.js, resulting in high latency and poor user experience on low-end mobile devices. This study presents a high-performance web virtualization architecture for AR, named TapTapp AR V11, which introduces a Nanite-style vision codec and a pure JavaScript execution pipeline. Using an applied research approach with a pretest-posttest experimental design, we evaluated the proposed system against the industry-standard MindAR framework. A sample of 45 compilation and tracking sessions was simulated under controlled conditions. The results demonstrate a dramatic improvement in performance: compilation time was reduced from approximately 23.50 seconds to 1.69 seconds, while the output file size decreased by 86\%. Statistical analysis using a paired Student's t-test confirmed the significance of these improvements (p $<$ 0.05). These findings validate the efficacy of eliminating heavy dependencies in favor of optimized binary processing and Fourier positional encoding, offering a scalable solution for ubiquitous web-based AR.
\end{abstract}
